\documentclass[../../main.tex]{subfiles}
\begin{document}
% SOURCE p=246; printed=231; object=ch07-products-exercises-title; region=middle
\Needspace{13\baselineskip}
\subsection*{习题}
% SOURCE p=246; printed=231; object=ch07-ex28-ordinal-product-borel; number=28; region=middle
\begin{exercise}\label{exercise:ch07-ordinal-product-borel}
若 $X$ 为所有不超过第一个不可数序数 $\omega_1$ 的序数组成的集合，赋予序拓扑，则 $\mathscr B_{X\times X}\ne\mathscr B_X\otimes\mathscr B_X$。事实上，$\{(x,y):x<y<\omega_1\}$ 是开集，但不属于 $\mathscr B_X\otimes\mathscr B_X$。（结合\cref{exercise:ch07-ordinal-measure}重新考察\cref{exercise:ch02-nonmeasurable-separately-measurable}。）
\end{exercise}
% SUPPLEMENT: id=v14-ch07-ex28-solution; parent=ch07-ex28-ordinal-product-borel; kind=solution+motivation+hypotheses; origin=independent-derivation
\begin{solution}
\textbf{完整补写（重构新增）。}记 $\Omega=\{x:x<\omega_1\}$，$E=\{(x,y):x<y<\omega_1\}$。先证 $E$ 开。对 $(x,y)\in E$，若存在 $z$ 满足 $x<z<y$，则开矩形 $[0,z)\times(z,\omega_1)$ 包含 $(x,y)$ 且包含于 $E$。若没有中间序数，则 $y$ 是 $x$ 的后继，开矩形 $[0,y)\times(x,\omega_1)$ 也具有这两个性质：其第一坐标至多 $x$，第二坐标至少 $y$。因此每点都有包含于 $E$ 的开矩形。

用\cref{exercise:ch07-ordinal-measure}中的 Borel 概率测度 $\lambda$：可数集质量为零，包含 $\Omega$ 中不可数闭集的 Borel 集质量为一，且 $\lambda(\{\omega_1\})=0$。对 $y<\omega_1$，$E^y=[0,y)$ 可数，所以 $\lambda(E^y)=0$；$y=\omega_1$ 时截面为空，也为零。对 $x<\omega_1$，$E_x=(x,\omega_1)$ 包含 $\Omega$ 中的闭尾段 $[x+1,\omega_1)$，该尾段不可数，故 $\lambda(E_x)=1$；$x=\omega_1$ 时为零。于是
\begin{align*}
\int_X\lambda(E^y)\,\mathrm d\lambda(y)=0,\qquad
\int_X\lambda(E_x)\,\mathrm d\lambda(x)=\lambda(\Omega)=1.
\end{align*}
若 $E$ 属于乘积 $\sigma$ 代数，因为 $\lambda$ 是有限测度，普通 Tonelli 定理要求这两个积分相等，矛盾。因此 $E$ 是所求的开 Borel 集，却不属于乘积 $\sigma$ 代数。
\end{solution}
\begin{remark}
\textbf{动机与条件（重构新增）。}利用全部初段可数而尾段不可数，使两个方向的截面质量分别为零和一，就能检测乘积可测性的失败。辅助测度 $\lambda$ 不 Radon，却是有限 Borel 测度，已足以应用普通 Tonelli；若误把它当作 Radon 测度，再用本节的 Radon 版本，会掩盖反例的关键。$X$ 虽紧 Hausdorff，缺少第二可数性，故不与\cref{theorem:ch07-product-borel-radon}矛盾。
\end{remark}
% END SUPPLEMENT: id=v14-ch07-ex28-solution
% SOURCE p=246; printed=231; object=ch07-ex29-large-discrete-product-borel; number=29; region=middle
\begin{exercise}\label{exercise:ch07-large-discrete-product-borel}
若 $X$ 为基数大于 $\mathfrak c$ 的集合，赋予离散拓扑，则 $\mathscr B_{X\times X}\ne\mathscr B_X\otimes\mathscr B_X$。事实上，对角线 $D=\{(x,y):x=y\}$ 闭，但不属于 $\mathscr B_X\otimes\mathscr B_X$。（使用\cref{exercise:ch01-countable-generators,proposition:ch01-transfinite-sigma-generation}。若 $D\in\mathscr B_X\otimes\mathscr B_X$，则 $D$ 属于由可数个矩形生成的某个子 $\sigma$ 代数 $\mathcal M$，进而 $D\in\mathcal N\otimes\mathcal N$，其中 $\mathcal N$ 是 $\mathscr B_X$ 的可数生成子 $\sigma$ 代数。于是所有 $x$ 都有 $\{x\}=D_x\in\mathcal N$，但 $\card(\mathcal N)\leqslant\mathfrak c$。）将 $X$ 换成其单点紧化，同样的论证仍成立。
\end{exercise}
% SUPPLEMENT: id=v14-ch07-ex29-solution; parent=ch07-ex29-large-discrete-product-borel; kind=solution+motivation+hypotheses; origin=independent-derivation
\begin{solution}
\textbf{完整补写（重构新增）。}离散 $X$ 的乘积仍离散，所以 $D$ 闭。若它在乘积 $\sigma$ 代数中，由\cref{exercise:ch01-countable-generators}，存在可数个矩形 $A_n\times B_n$，使 $D\in\sigma(\{A_n\times B_n:n\geqslant1\})$。令 $\mathcal N=\sigma(\{A_n,B_n:n\geqslant1\})$，则全部这些矩形属于 $\mathcal N\otimes\mathcal N$，所以 $D$ 也属于其中。固定 $x$ 后取截面的运算保持补和可数并；生成矩形的截面等于某个 $B_n$ 或空集。因此每个 $\mathcal N\otimes\mathcal N$ 中集合的截面属于 $\mathcal N$，特别地 $\{x\}=D_x\in\mathcal N$。

可数生成的 $\sigma$ 代数至多有 $\mathfrak c$ 个元素。为明确这个基数估计，用\cref{proposition:ch01-transfinite-sigma-generation}按所有可数序数分层生成：初始层可数，若一层基数至多 $\mathfrak c$，取补不增基数，所有可数序列组成的集合基数至多 $\mathfrak c^{\aleph_0}=\mathfrak c$，故取可数并后仍至多 $\mathfrak c$；可数极限层也满足此界。所有可数序数的基数 $\aleph_1\leqslant\mathfrak c$，因此全部层的并至多有 $\mathfrak c$ 个元素。但不同的 $x$ 给不同的单点集，迫使 $\card(\mathcal N)\geqslant\card(X)>\mathfrak c$，矛盾。

若换成单点紧化 $X^*=X\cup\{\infty\}$，它是紧 Hausdorff 空间，故其对角线闭：若 $u\ne v$，可取不交开邻域，乘积邻域与对角线不交，从而对角线补集开。若对角线在 $\mathscr B_{X^*}\otimes\mathscr B_{X^*}$ 内，重复可数矩形生成及截面论证，得到一个含全部单点集的可数生成 $\sigma$ 代数，仍与 $\card(X^*)>\mathfrak c$ 矛盾。
\end{solution}
\begin{remark}
\textbf{动机与条件（重构新增）。}一个具体的乘积可测集至多动用了可数个矩形生成元，因此其截面种类被一个至多连续统大小的 $\sigma$ 代数约束；对角线却要求每个点产生不同截面。严格大于 $\mathfrak c$ 是此基数矛盾的条件，不能把本证明直接改成“不可数即可”。单点紧化说明即使整体紧，也不能自动令 Borel 乘积与乘积 Borel 相同。
\end{remark}
% END SUPPLEMENT: id=v14-ch07-ex29-solution
% SOURCE p=246; printed=231; object=ch07-ex30-lsc-product-integrals; number=30; region=lower
\begin{exercise}\label{exercise:ch07-lsc-product-integrals}
设 $\mu,\nu$ 分别为 $X,Y$ 上的 Radon 测度，不必 $\sigma$ 有限。若 $f$ 是 $X\times Y$ 上非负下半连续函数，则 $x\mapsto\int f_x\,\mathrm d\nu$、$y\mapsto\int f^y\,\mathrm d\mu$ Borel 可测，且
\begin{align*}
\int f\,\mathrm d(\mu\mathbin{\widehat\times}\nu)
=\iint f\,\mathrm d\mu\,\mathrm d\nu
=\iint f\,\mathrm d\nu\,\mathrm d\mu.
\end{align*}
\end{exercise}
% SUPPLEMENT: id=v14-ch07-ex30-solution; parent=ch07-ex30-lsc-product-integrals; kind=solution+motivation+hypotheses; origin=independent-derivation
\begin{solution}
\textbf{完整补写（重构新增）。}令 $\mathcal G=\{g\in C_c(X\times Y):0\leqslant g\leqslant f\}$。它非空，因为包含零函数；且对有限最大值封闭，所以向上有向。由\cref{proposition:ch07-lsc-properties}，$f=\sup_{g\in\mathcal G}g$，故固定坐标后的截面也分别为 $f_x=\sup g_x$、$f^y=\sup g^y$。这些截面本身下半连续，因为是下半连续函数与连续嵌入的复合。

对每个固定 $x$，由有向下半连续收敛定理，
\begin{align*}
H(x):=\int f_x\,\mathrm d\nu
=\sup_{g\in\mathcal G}\int g_x\,\mathrm d\nu.
\end{align*}
右侧为连续函数 $H_g(x)=\int g_x\,\mathrm d\nu$ 的上确界，因此 $H$ 下半连续，特别 Borel。$H_g$ 族向上有向，因为 $H_{\max(g,h)}\geqslant H_g,H_h$。再次使用同一定理及连续紧支撑积分公式，得
\begin{align*}
\int H\,\mathrm d\mu
&=\sup_{g\in\mathcal G}\int H_g\,\mathrm d\mu
=\sup_{g\in\mathcal G}\int g\,\mathrm d(\mu\mathbin{\widehat\times}\nu)\\
&=\int f\,\mathrm d(\mu\mathbin{\widehat\times}\nu).
\end{align*}
最后一项是在乘积空间上再用有向下半连续收敛定理。交换 $X,Y$，$J(y)=\int f^y\,\mathrm d\mu$ 为连续部分积分的上确界，故下半连续；完全相同的有向上确界计算给 $\int J\,\mathrm d\nu=\int f\,\mathrm d(\mu\mathbin{\widehat\times}\nu)$，完成两种积分次序。
\end{solution}
\begin{remark}
\textbf{动机与条件（重构新增）。}下半连续性使函数可由连续紧支撑函数的有向上确界恢复，从而可用拓扑正则性取代全局 $\sigma$ 有限性。非负性保证所有积分及上确界都合法，即使为 $+\infty$ 也不相减。只要求任意 Borel 可测就不能使用这套有向逼近论证，前面的非 $\sigma$ 有限反例说明不能直接照搬更强结论。
\end{remark}
% END SUPPLEMENT: id=v14-ch07-ex30-solution
% SOURCE p=246; printed=231; object=ch07-ex31-product-baire; number=31; region=lower
\begin{exercise}\label{exercise:ch07-product-baire}
关于乘积空间中的 Baire 集，有如下结论。
\begin{enumerate}
\item\label{item:ch07-product-baire-inclusion} $\mathscr B_{X\times Y}^0\subseteq\mathscr B_X^0\otimes\mathscr B_Y^0$。（提示：\cref{proposition:ch07-cc-product-integrals}的可测性结论把 $\mathscr B$ 换成 $\mathscr B^0$ 后仍成立。）
\item\label{item:ch07-product-baire-equality} 若 $X,Y$ 均为紧空间或第二可数空间，则 $\mathscr B_{X\times Y}^0=\mathscr B_X^0\otimes\mathscr B_Y^0$。
\item\label{item:ch07-discrete-product-baire-strict} 若 $X$ 为不可数离散空间，则 $\mathscr B_{X\times X}^0\ne\mathscr B_X^0\otimes\mathscr B_X^0$。
\end{enumerate}
\end{exercise}
% SUPPLEMENT: id=v14-ch07-ex31-solution; parent=ch07-ex31-product-baire; kind=solution+motivation+hypotheses; origin=independent-derivation
\begin{solution}
\textbf{逐项补写（重构新增）。}第一项：每个 $g\in C_c(X)$、$h\in C_c(Y)$ 分别是两因子的 Baire 可测函数，所以 $(g\circ\pi_X)(h\circ\pi_Y)$ 对 $\mathscr B_X^0\otimes\mathscr B_Y^0$ 可测。有限和及逐点极限仍可测，张量逼近定理使每个 $C_c(X\times Y)$ 函数都对这个乘积代数可测。$\mathscr B_{X\times Y}^0$ 是使所有这类函数可测的最小 $\sigma$ 代数，因此得到所求包含关系。

第二项：先说明每个因子都存在一列 $C_c$ 函数 $u_n$，使 $u_n\to1$ 逐点。紧空间直接取 $u_n=1$；第二可数 LCH 空间有可数个相对紧开集覆盖，取其前 $n$ 个闭包的并为 $K_n$，由局部 Urysohn 引理选择 $0\leqslant u_n\leqslant1$、$u_n\in C_c$ 且 $u_n=1$ 于 $K_n$，则每点最终值为 $1$。两个因子中分别记此序列为 $u_n,v_n$。

给定 $g\in C_c(X)$，函数 $g(x)v_n(y)$ 属于 $C_c(X\times Y)$，所以为乘积空间的 Baire 函数；它逐点趋于 $g\circ\pi_X$，故该复合函数也 Baire。由 $\mathscr B_X^0$ 的生成定义，这说明投影 $\pi_X:(X\times Y,\mathscr B_{X\times Y}^0)\to(X,\mathscr B_X^0)$ 可测。具体地，所有在投影下原像为 Baire 集的 $X$ 子集组成 $\sigma$ 代数，且包含每个 $g\in C_c(X)$ 的 Borel 原像，因而包含 $\mathscr B_X^0$。对 $h\in C_c(Y)$ 用 $u_n(x)h(y)$ 同理证明另一投影可测。因此全部 Baire 矩形在 $\mathscr B_{X\times Y}^0$ 内，得到反向包含，与第一项合并即为等式。这个证明还覆盖一个因子紧、另一个因子第二可数的混合情形。

第三项：不可数离散空间 $X$ 上 $C_c$ 函数恰为有限支撑函数，它们都对可数或余可数 $\sigma$ 代数可测，故 $\mathscr B_X^0$ 包含于这个代数。反向，每个单点示性函数都在 $C_c$，所以所有单点集为 Baire 集，其可数并和补集给出全部可数或余可数集。于是 $\mathscr B_X^0$ 正是该代数。离散乘积 $X\times X$ 同理，其 Baire 代数也由可数或余可数集合组成。

固定 $x_0\in X$，矩形 $\{x_0\}\times X$ 属于 $\mathscr B_X^0\otimes\mathscr B_X^0$，但它不可数，其补集 $(X\setminus\{x_0\})\times X$ 也不可数，故不属于 $\mathscr B_{X\times X}^0$。结合第一项，包含关系在此严格。
\end{solution}
\begin{remark}
\textbf{动机与条件（重构新增）。}Baire 结构由紧支撑连续函数生成，所以其乘积包含方向与 Borel 情形相反。要恢复投影的可测性，需要把另一个变量上的常数 $1$ 作为紧支撑函数的逐点序列极限；紧性或第二可数的 LCH 结构正好提供这样的列。不可数离散空间的紧支集只能有限，任何一列这类函数只能涉及可数多个点，因此这一恢复办法失效，题中的严格包含反例便说明其后果。
\end{remark}
% END SUPPLEMENT: id=v14-ch07-ex31-solution
\end{document}
